\scanpage{39}
\mspage{13}

[NB \deleted{\unsure{Pour}}On pense notamment au cas où $X$ est associé à une \uline{carte}, et $H$ est le groupe des autom. de cette carte, qui est fini -- la condition $X^H=\varnothing$ signifie aussi qu'il n'y a pas de facette de la carte fixée par $H$...]

Dans tous les cas où $Q_X$ est défini, $Q_X$ est hyperbolique (i.e. : courbure $<0$) sss $X$ est hyperbolique, et alors $Q_X$ est de courbure \inserted{constante} égale à $-1$. D'autre part, $Q_X$ est \uline{plate} \deleted{dans les cas} (ou encore, à courbure nulle) sss $\widetilde X\simeq\mathbb C$ [y inclus le cas a) ci-dessus, où on a encore une $Q_X$ mod cte]. Le cas \unsure{restant|restreint} est celui du plan projectif réel, conforme, \deleted{qui est} \deleted{\unsure{dans le}}\inserted{pour} lequel $Q_X$ est une métrique homogène, de courbure Cte \deleted{égale à} et $>0$.

Donc chaque fois que $Q_X$ est défini, ou défini mod cte, il est à courbure constante, \deleted{\unsure{et est}} Il en est encore de = de $Q_{X,H}$ quand $X\simeq\mathbb P^1_{\mathbb C}$ et $H\subset\mathrm{Aut}_{\mathrm{conf}}X$ est cpt et tel que $X^H=\varnothing$, on trouve \deleted{une} métriques sph. euclidiennes, \unsure{donc|dans} de courbure constante $>0$.

% slanted
\margin{\uline{NB} Pour tte $X$ 0-connexe conforme, on \deleted{\illegible{1}} a dans l'ensemble \deleted{des métriques} $\mathrm{Riem}(X)$ des métr. riemanniennes comp. avec la str. conforme une orbite canonique sous $\mathrm{Aut}_{\mathrm{conf}}(X)$, qui est réduit à 1 pt dans tous les cas, sauf les cas a) et b) de la p. 12, où dans le cas a), on trouve une demi-droite dans $\mathrm{Riem}(X)$, et dans le cas b), une variété de dim 3 (\deleted{\unsure{qui|quelle}}\inserted{telle} \deleted{\unsure{les}} que $Q,Q'\in R(X)$ $Q\ne Q'$ implique $Q'$ et $Q$ \uline{non} prop.) Notion de métriques riemanniennes distinguées (a un sens aussi dans le cas \unsure{disconnexe|discontinu},}
% archive sheet number: 41
